\documentclass[a4paper,12pt]{article}
\usepackage[utf8]{inputenc}
\usepackage[T1]{fontenc}
\usepackage[italian]{babel}
\usepackage[a4paper,margin=2.5cm]{geometry}
\usepackage{setspace}
\onehalfspacing

\title{Sicurezza nelle console di gioco}
\author{Marica Bottega, Enrico Cicciarella}
\date{\today}

\begin{document}
\maketitle
\thispagestyle{empty}

%% TODO: Inserire indice

\paragraph{Introduzione}
Lo scopo di questo elaborato è quello di analizzare come le tecniche offensive e difensive nelle console di gioco casalinghe si siano evoluti nel tempo.\\
Il modello di minaccia all'interno del quale sono immersi questi sistemi è il \textbf{modello MATE} -- \textit{Man At The End}, ovvero quello in cui l'attaccante ha completa disposizione fisica della macchina.\\
Analizziamo quindi lo stato dell'arte delle tecniche di difesa impiegati dalle più grandi case produttrici di console di gioco, e di come siano comunque state prese di mira attraverso exploit software, corruzione della memoria e attacchi hardware come la \textit{Fault Injection}, con l’obiettivo di aggirare la \textit{Chain of Trust}.

\paragraph{Memory Corruption}
\paragraph{Crittografia White-Box}
Il contesto storico è quello dei primi anni 2000 -- periodo in cui erano appena entrate in commercio console come xBox, GameCube e PlayStation 2. \\
Come anticipato in introduzione, il modello \textit{MATE} assume in partenza che tutto ciò che si trova all'esterno del processore (RAM, bus di sistema, periferiche, \ldots) sia esposto e potenzialmente sotto il controllo dell'avversario.\\
Il principio è lo stesso che giustifica la crittografia White-Box: chi esegue il programma è lo stesso che può attaccarlo. \\
Il contributo di \textit{Weidong Shi, Hsien-Hsin S. Lee, Chenghuai Lu e Tao Zhang}, che esamina \textbf{Attacchi e analisi del rischio per i sistemi hardware a supporto della protezione dalla copia del software} cita .
\paragraph{Fault Injection}
\paragraph{Architetture Moderne}
\paragraph{Bibliografia}


\end{document}
